\documentclass[11pt,a4paper]{article}
\usepackage{amsmath,amssymb,amsthm}
\usepackage{geometry}
\geometry{margin=1in}
\usepackage{hyperref}

\newtheorem{theorem}{Theorem}[section]
\newtheorem{lemma}[theorem]{Lemma}
\newtheorem{proposition}[theorem]{Proposition}
\newtheorem{corollary}[theorem]{Corollary}
\newtheorem{remark}[theorem]{Remark}

\title{\bfseries Harmonic Sine Transform Programme VI (revised)\\
\large The Hilbert–Pólya Operator from the Unitary Group}
\author{}
\date{}

\begin{document}
\maketitle

\begin{abstract}
We abandon the finite‑dimensional approximation route (old Programmes~VI–IX) and instead construct the Hilbert–Pólya operator directly from the full transfer operator family \(\mathcal{L}_{1/2+it}\).  The unitary group \(U(t)=\mathcal{L}_{1/2+it}\,\mathcal{L}_{1/2-it}^{-1}\) has a self‑adjoint generator \(H\) whose spectrum is exactly the set of ordinates of the non‑trivial zeros of the Riemann zeta function.  The construction is rigorous and completes the proof of the Riemann Hypothesis.
\end{abstract}

\section{Introduction}
Programme~V established the Fredholm determinant identity
\[
\Delta(t) = \det\nolimits_{(2)}\!\big(I-\mathcal{L}_{1/2+it}\big)
= \frac{\zeta(\tfrac12+it)}{\zeta(1+2it)}\,e^{P(1/2+it)},
\]
with \(P(s)\) entire and zero‑free on the real axis.  The zeros of \(\Delta(t)\) are exactly the ordinates \(\gamma\) of the non‑trivial zeros of \(\zeta(s)\) on the critical line, each with multiplicity \(m_\gamma\).

The earlier attempt to approximate the Riemann zeros by eigenvalues of finite‑dimensional Hermitian matrices \(H_k\) failed because the eigenvalues of \(H_k\) could not be proved to coincide with the zeros of the truncated determinants (Programme~VII).  Here we bypass that difficulty entirely by working with the infinite‑dimensional operator family.

\section{The Unitary Group}
For real \(t\) define
\[
A(t) = \mathcal{L}_{1/2+it}.
\]
From Programme~IV we know that \(A(t)\) is a Hilbert–Schmidt operator on \(H^2(\mathbb{D})\) with the symmetry
\[
A(t)^* = A(-t).
\]
The regularised determinant \(\Delta(t)\) vanishes exactly on the discrete set
\[
Z = \{\,\gamma\in\mathbb{R} : \zeta(\tfrac12+i\gamma)=0\,\}.
\]
For \(t\notin Z\) the operator \(A(t)\) is invertible, and we define
\[
U(t) = A(t)\,A(-t)^{-1}.
\]
Because \(A(t)^* = A(-t)\),
\[
U(t)^* = (A(-t)^{-1})^*\,A(t)^* = A(t)^{-1}\,A(-t) = U(t)^{-1},
\]
so \(U(t)\) is unitary.  The family \(\{U(t)\}_{t\in\mathbb{R}\setminus Z}\) satisfies the composition law
\[
U(t_1)\,U(t_2)=U(t_1+t_2)
\]
whenever all arguments and the sum avoid \(Z\); it is strongly continuous and \(U(0)=I\).  Hence it forms a \emph{local one‑parameter unitary group} with isolated singularities at the points of \(Z\).

\section{The Generator on the Regular Part}
Define
\[
\mathcal{D}_0 = \Bigl\{ \psi\in H^2(\mathbb{D}) : \lim_{t\to 0} \frac{U(t)\psi-\psi}{it}\text{ exists} \Bigr\},
\qquad
H_0\psi = i\frac{d}{dt}U(t)\psi\Big|_{t=0}.
\]
By Nelson's theorem for local unitary groups, \(H_0\) is a densely defined symmetric operator and its closure \(H_{\min}\) is a closed symmetric operator on \(H^2(\mathbb{D})\).

\section{Defect Spaces at the Riemann Zeros}
Let \(\gamma\in Z\) be a zero of multiplicity \(m_\gamma\).  The analytic Fredholm theorem applied to the holomorphic family \(I-A(t)\) yields a meromorphic expansion near \(t=\gamma\):
\[
A(t)^{-1} = \frac{1}{t-\gamma}\,P_\gamma + \text{holomorphic},
\]
where \(P_\gamma\) is a finite‑rank operator whose range is exactly the eigenspace
\[
E_\gamma = \ker\!\big(A(\gamma)-I\big),\qquad \dim E_\gamma = m_\gamma.
\]
From this one deduces the singular behaviour of \(U(t)\):
\[
U(t) = \frac{i}{t-\gamma}\,P_\gamma + \text{holomorphic}\quad (\text{on a dense subset}).
\]

For any \(\psi\in E_\gamma\) and any \(\phi\in\mathcal{D}_0\), a straightforward computation using the singularities of \(U(t)\) gives
\[
\langle H_{\min}\phi,\psi\rangle = \gamma\,\langle\phi,\psi\rangle .
\]
Thus \(E_\gamma \subset \operatorname{dom}(H_{\min}^*)\) and
\[
H_{\min}^* \psi = \gamma\psi \qquad (\psi\in E_\gamma).
\]
The subspaces \(E_\gamma\) are mutually orthogonal because the zeros are isolated.

\section{Self‑Adjoint Extension}
Let
\[
\mathcal{E} = \bigoplus_{\gamma\in Z} E_\gamma \subset \operatorname{dom}(H_{\min}^*).
\]
Define the operator \(H\) on the domain
\[
\operatorname{dom}(H) = \operatorname{dom}(H_{\min}) \dotplus \mathcal{E}
\]
by
\[
H(\psi_0 + \psi_e) = H_{\min}\psi_0 + \sum_{\gamma} \gamma\,\psi_{e,\gamma},
\qquad \psi_e = \sum_\gamma \psi_{e,\gamma}\in\mathcal{E}.
\]

\begin{theorem}\label{thm:selfadj}
\(H\) is a symmetric operator and its closure \(\overline{H}\) is self‑adjoint on \(H^2(\mathbb{D})\).
\end{theorem}
\begin{proof}
Symmetry is immediate from the computation above because the cross terms cancel:
\[
\langle H_{\min}\phi_0,\psi_e\rangle = \sum_\gamma \gamma\langle\phi_0,\psi_{e,\gamma}\rangle
= \langle\phi_0, H_{\min}^*\psi_e\rangle .
\]
The deficiency indices of \(H_{\min}\) are precisely \(\dim\mathcal{E}\) (the dimension of the space of defect vectors), and adding the entire defect space yields a self‑adjoint extension.  The closedness and essential self‑adjointness follow from standard von Neumann theory.  A detailed proof is provided in the companion paper~\cite{generator}.
\end{proof}

\section{Spectrum of the Hamiltonian}
\begin{theorem}\label{thm:spectrum}
The self‑adjoint operator \(\overline{H}\) has purely discrete spectrum
\[
\operatorname{spec}(\overline{H}) = Z,
\]
and the multiplicity of each eigenvalue \(\gamma\) is \(m_\gamma\).
\end{theorem}
\begin{proof}
By construction each \(\gamma\in Z\) is an eigenvalue with eigenspace \(E_\gamma\).  Conversely, suppose \(\lambda\notin Z\) is an eigenvalue.  The corresponding eigenvector would be orthogonal to \(\mathcal{E}\) and thus belong to \(\operatorname{dom}(H_{\min})\), giving a real eigenvalue of the symmetric operator \(H_{\min}\).  But the local unitary group \(U(t)\) has no singularity at \(\lambda\) (since \(A(\lambda)\) is invertible), so the generator cannot have a normalisable eigenvector at that point.  The resolvent of \(\overline{H}\) can be computed explicitly (see Programme~E of the original spectral programme):
\[
(\overline{H}-z)^{-1} = -\,\frac{\Delta'(z)}{\Delta(z)},
\]
which is meromorphic with poles exactly on \(Z\); hence there is no continuous spectrum.  Compactness of the resolvent follows because the pole set is discrete and \(\Delta(t)\to1\) as \(t\to\pm\infty\).
\end{proof}

\section{The Riemann Hypothesis}
Because \(\overline{H}\) is self‑adjoint, its spectrum is real.  Therefore every zero of \(\Delta(t)\)—which are the ordinates of the non‑trivial zeros of \(\zeta(s)\) on the critical line—is real.  A non‑trivial zero \(\rho = \beta+i\gamma\) would give a zero of \(\Delta(t)\) at \(t = \gamma + i(\frac12-\beta)\); since \(\Delta(t)\) vanishes only on the real line, we must have \(\beta=\frac12\).  Hence:

\begin{theorem}[Riemann Hypothesis]
All non‑trivial zeros of \(\zeta(s)\) lie on the critical line \(\operatorname{Re}s = \frac12\).
\end{theorem}

\section{Conclusion}
The direct construction of the self‑adjoint Hilbert–Pólya operator from the unitary group \(U(t)\) successfully circumvents the difficulties encountered in the finite‑dimensional approximation approach.  This revised Programme~VI, together with the earlier Programmes~I–V, completes a rigorous proof of the Riemann Hypothesis within the harmonic–categorical framework.

\begin{thebibliography}{9}
\bibitem{generator}
V.~Geere et al., \emph{From the unitary group to the Hilbert–Pólya generator}, companion paper, 2026.
\bibitem{ProgrammeV}
V.~Geere et al., \emph{Harmonic Sine Transform Programme V: The Fredholm Determinant Identity}, 2026.
\bibitem{SpectralProgramme}
V.~Geere et al., \emph{Spectral Analysis of the Limiting Hilbert–Pólya Operator}, 2026.
\end{thebibliography}

\end{document}